\documentclass[11pt]{article}

\usepackage[T1]{fontenc}
\usepackage[utf8]{inputenc}
\usepackage[a4paper,margin=1in]{geometry}
\usepackage{lmodern}
\usepackage{microtype}
\usepackage{amsmath,amssymb,amsthm}
\usepackage{graphicx}
\usepackage{booktabs}
\usepackage{array}
\usepackage{enumitem}
\usepackage{caption}
\usepackage{hyperref}
\hypersetup{colorlinks=true,linkcolor=blue,citecolor=blue,urlcolor=blue}

\newtheorem{proposition}{Proposition}
\newtheorem{remark}{Remark}
\newtheorem{definition}{Definition}
\DeclareMathOperator*{\argmin}{arg\,min}

\title{Paper 001: Chart-Scaled Cubic Pandrosion by Finite-Slope Halley Acceleration\\
\large Thales scaling, Riemann inversion, and a rebuilt benchmark}
\author{Ivan Besevic\\Independent Mathematical Research\\\texttt{github.com/ivan-fr/poussiere\_de\_x}}
\date{May 2026}

\begin{document}
\maketitle

\begin{abstract}
Paper 000 developed the monomial Pandrosion map for $z^p=x$ and showed that its residual multiplier is controlled by the normalized target.  The practical rule was geometric: choose a homothety, and when useful a reciprocal Riemann chart, so that the normalized target $y$ is close to $1$ before iterating.  The previous 001 note introduced a derivative-free cubic finite-slope Halley correction built from Pandrosion probes.  This rebuilt version combines the two ideas.

For the monomial equation $z^p=x$, we first choose a chart and a root-scale $B$.  In the direct chart we solve $u^p=x/B^p$ and return $z=Bu$.  In the reciprocal chart we solve $u^p=1/(xB^p)$ and return $z=1/(Bu)$.  The finite-slope Halley-Pandrosion step is then applied only to the normalized equation $u^p=y$, starting from $u_0=1$.  On a scale palette $B\in\{q^{m+j/K}:m\in\mathbb Z,\,j=0,\ldots,K-1\}$, the same logarithmic covering argument from Paper 000 gives $1\le y\le q^{p/(2K)}$, hence $1\le y^{1/p}\le q^{1/(2K)}$.  This moves the finite-slope problem into a magnitude-free root-coordinate neighborhood.

The 001 benchmark has also been rebuilt.  The original local polynomial benchmark still shows the same hierarchy: raw finite-slope Pandrosion is fragile at high degree, Steffensen-Pandrosion is a strong quadratic upgrade, and cubic finite-slope Halley remains the most robust method on the tested local families.  A new chart-scaling benchmark over 240 decades of monomial targets shows the extra effect of Thales/Riemann preconditioning: with $q=10$ and $K=4$, cubic H3 reaches 100\% success for all tested degrees $p\in\{3,5,8,16,32,64\}$ within 25 updates, while the unscaled chart succeeds only on a small fraction of the same grid.
\end{abstract}

\paragraph{One-sentence summary.}
The cubic finite-slope Halley-Pandrosion correction becomes a chart-scaled monomial solver by first using Thales homothety and Riemann inversion to make $|\log y|$ small, then applying the same derivative-free cubic update to the normalized equation $u^p=y$.

\section{Introduction}
The original 001 note asked whether the finite-slope/Pandrosion philosophy could move beyond Steffensen's quadratic acceleration without using exact derivatives.  The answer was the cubic update
\[
H_3(a)=a-\frac{2f(a)D_1(a)}{2D_1(a)^2-f(a)D_2(a)},
\]
where $D_1$ and $D_2$ are symmetric finite-slope approximants of $f'$ and $f''$.

This revision adds the missing geometric preconditioner from Paper 000.  Paper 000 showed that, for monomial Pandrosion, the contraction fingerprint depends on the normalized target $y$, not on the original magnitude of $x$.  The correct geometric action is therefore not to push one affine chart across all scales, but to choose a chart and homothety first.  The same principle is useful for the finite-slope cubic method: the finite-slope probes are reliable when the root-coordinate problem is local, and chart scaling makes it local before the probes are formed.

The aim is still modest.  Newton, Halley, and standard Householder methods remain the natural production tools for monomial root extraction when derivatives or special formulas are allowed.  The point here is different: keep the finite-slope derivative-free structure, apply the Thales/Riemann scaling rule, and measure the effect on the 001 solver ladder.

\section{The finite-slope solver ladder}
Let $f:\mathbb R\to\mathbb R$ have a simple root $r$, with $f(r)=0$ and $f'(r)\ne0$.  Given two points $a$ and $b$, write
\[
f[a,b]=\frac{f(b)-f(a)}{b-a}.
\]
The numerical parameters used throughout the rebuilt benchmark are
\[
\rho=0.10,\qquad \eta=0.75.
\]

\subsection{Raw finite-slope Pandrosion}
The raw step uses the homothetic probe
\[
h_0(a)=\rho(1+|a|)
\]
and defines
\[
P(a)=a-\frac{f(a)}{f[a,a+h_0(a)]}.
\]
This is the scalar finite-slope primitive.  It is usually only first-order accurate as an iteration, but it gives the predictor scale used below.

\subsection{Steffensen-Pandrosion}
Aitken-Steffensen acceleration applied to the raw map $P$ gives
\[
S(a)=a-\frac{(P(a)-a)^2}{P(P(a))-2P(a)+a},
\]
when the denominator is non-zero.  This cancels the leading linear term of the raw fixed-point error and is locally quadratic under the usual non-degeneracy assumptions.

\subsection{Cubic finite-slope Halley-Pandrosion}
First compute the raw predictor displacement
\[
\delta_{\rm raw}(a)=-\frac{f(a)}{f[a,a+h_0(a)]}.
\]
Then choose the symmetric local probe radius
\[
h(a)=\eta\min\bigl(|\delta_{\rm raw}(a)|,h_0(a)\bigr),
\]
with a tiny floating-point floor in the implementation.  Define
\[
D_1(a)=\frac{f(a+h(a))-f(a-h(a))}{2h(a)},
\qquad
D_2(a)=\frac{f(a+h(a))-2f(a)+f(a-h(a))}{h(a)^2}.
\]
The cubic finite-slope Halley-Pandrosion update is
\begin{equation}
H_3(a)=a-\frac{2f(a)D_1(a)}{2D_1(a)^2-f(a)D_2(a)}.
\label{eq:H3}
\end{equation}
It is exactly Halley's rational correction with $f'$ and $f''$ replaced by symmetric finite-slope data.

\begin{proposition}[formal local cubic behavior]
Assume $f\in C^4$ near a simple root $r$, write $e=a-r$, and suppose the adaptive probe satisfies
\[
h(a)=\theta e+O(e^2)
\]
for some finite non-zero constant $\theta$.  Then the finite-slope Halley-Pandrosion update satisfies
\[
H_3(a)-r=C_\theta e^3+O(e^4)
\]
for a finite constant $C_\theta$ depending on $f$ and on $\theta$.
\end{proposition}

\begin{proof}[Proof sketch]
The symmetric formulas give
\[
D_1(a)=f'(a)+\frac{h(a)^2}{6}f'''(a)+O(h(a)^4),
\qquad
D_2(a)=f''(a)+O(h(a)^2)+O(eh(a)).
\]
When $h(a)=O(e)$, this gives $D_1(a)=f'(a)+O(e^2)$ and $D_2(a)=f''(a)+O(e)$.  Substitution into Halley's rational expression preserves the classical cancellation through the quadratic error term, so the remaining defect is cubic.
\end{proof}

The raw predictor is therefore not cosmetic.  It supplies a derivative-free local scale estimate.  The chart-scaling layer below makes that estimate reliable over very different magnitudes of $x$.

\section{Thales scaling and Riemann inversion for the cubic step}
Consider the monomial root problem
\[
z^p=x,\qquad x>0,\quad p\ge2.
\]
A root-coordinate homothety and a reciprocal chart produce two normalized equations.

\paragraph{Direct Thales chart.}
Set $z=Bu$ with $B>0$.  Then
\begin{equation}
u^p=y_+(B),\qquad y_+(B)=\frac{x}{B^p},\qquad z=Bu.
\label{eq:direct}
\end{equation}

\paragraph{Reciprocal Riemann chart.}
Set $z=1/(Bu)$.  Then
\begin{equation}
u^p=y_\infty(B),\qquad y_\infty(B)=\frac{1}{xB^p},\qquad z=\frac{1}{Bu}.
\label{eq:reciprocal}
\end{equation}
These are not different solvers.  They are the same normalized finite-slope solver transported through different charts.

\begin{proposition}[chart-scaled correctness]
If $u^p=x/B^p$, then $(Bu)^p=x$.  If $u^p=1/(xB^p)$, then $(1/(Bu))^p=x$.
\end{proposition}

\begin{proof}
The first identity is $B^pu^p=B^p(x/B^p)=x$.  The second is $(Bu)^{-p}=1/(B^pu^p)=1/(B^p/(xB^p))=x$.
\end{proof}

For the cubic finite-slope method, the normalized equation is
\[
f_y(u)=u^p-y,
\]
and the default normalized start is
\[
u_0=1.
\]
The chart layer is responsible for ensuring that the true normalized root $u_*=y^{1/p}$ is close to this start.

\section{Logarithmic palette rule}
Let
\[
\mathcal B_{q,K}=\{q^{m+j/K}:m\in\mathbb Z,\;j=0,\ldots,K-1\},\qquad q>1.
\]
For each $B\in\mathcal B_{q,K}$ and each chart $\chi\in\{+,\infty\}$, form $y_\chi(B)$ from \eqref{eq:direct} or \eqref{eq:reciprocal}.  The selection rule is
\begin{equation}
(\chi_*,B_*)\in\argmin\bigl\{\log y_\chi(B):B\in\mathcal B_{q,K},\;\chi\in\{+,\infty\},\;y_\chi(B)\ge1\bigr\}.
\label{eq:selector}
\end{equation}
The admissibility condition $y\ge1$ keeps the normalized problem on the same positive branch used in the monomial multiplier analysis.

The scale-palette theorem from Paper 000 can be restated in root-coordinate language as follows.

\begin{proposition}[root-coordinate bound]
With the selector \eqref{eq:selector}, every target $x>0$ admits a normalized target satisfying
\[
1\le y_*\le q^{p/(2K)}.
\]
Consequently the normalized root obeys
\begin{equation}
1\le u_*=y_*^{1/p}\le q^{1/(2K)}.
\label{eq:rootbound}
\end{equation}
\end{proposition}

\begin{proof}
On the logarithmic circle of circumference $p\log q$, $K$ equally spaced phase offsets have covering radius $p\log q/(2K)$.  Direct and reciprocal charts identify complementary positions on this circle.  Thus the selected normalized target has $0\le\log y_*\le p\log q/(2K)$.  Taking the $p$-th root gives \eqref{eq:rootbound}.
\end{proof}

For the benchmark choice $q=10$ and $K=4$, this gives
\[
1\le u_*\le 10^{1/8}\approx1.3335,
\]
independently of the original magnitude of $x$ and independently of $p$.  This is the mechanism by which the Thales/Riemann layer feeds the finite-slope cubic method: the cubic probes are formed after the problem has been moved to a uniformly local root-coordinate interval.

\section{Chart-scaled H3 algorithm}
The combined monomial algorithm is:
\begin{enumerate}[leftmargin=2.2em]
\item Choose $q,K$ and the scale palette $\mathcal B_{q,K}$.
\item Use \eqref{eq:selector} to choose a chart $\chi\in\{+,\infty\}$, a scale $B$, and a normalized target $y\ge1$.
\item Solve $f_y(u)=u^p-y=0$ from $u_0=1$ using raw Pandrosion, Steffensen-Pandrosion, or the cubic H3 update \eqref{eq:H3}.
\item Return $z=Bu$ in the direct chart or $z=1/(Bu)$ in the reciprocal chart.
\end{enumerate}

\begin{remark}[why inversion still matters]
For some small targets, direct scaling by a small $B$ is enough.  For others, the reciprocal chart gives the smaller normalized logarithmic gap.  The selector does not force either geometry; it lets the scale lattice decide.  In the examples below, $p=3$, $x=2\cdot10^{-6}$ is best normalized by a direct small scale, while $p=5$, $x=5\cdot10^{-6}$ is best normalized by the reciprocal chart with $B=10$.
\end{remark}

\section{Rebuilt local benchmark from 001}
The first benchmark reproduces the local scalar test from 001.  All families have the simple root $r=1$:
\[
f_1(t)=t^d-1,
\qquad
f_2(t)=t^d-t,
\qquad
f_3(t)=(t-1)(t^{d-1}+1/2).
\]
The degrees are $d\in\{2,4,8,16,32,64\}$, and the 201 starting points are uniformly spaced in $[0.9,1.1]$.  A run succeeds if
\[
|t_k-1|<10^{-10}\quad\text{or}\quad |f(t_k)|<10^{-10}
\]
within 25 updates.  Function evaluations are counted without caching; this makes one raw update cost three evaluations after the initial check, and one Steffensen or H3 update cost five evaluations after the initial check.

\begin{table}[htbp]
\centering
\caption{Rebuilt aggregate local benchmark, aggregated across the three polynomial families.}
\small
\input{tables/table_local_summary.tex}
\label{tab:local-summary}
\end{table}

\begin{figure}[htbp]
\centering
\includegraphics[width=0.96\linewidth]{assets/fig_convergence_trace.pdf}
\caption{Representative trace on $f(u)=u^{16}-1$ from $u_0=1.08$.  The curves separate into linear/raw, quadratic/Steffensen, and cubic/H3 behavior until they hit the plotted precision floor.}
\label{fig:trace}
\end{figure}

\begin{figure}[htbp]
\centering
\includegraphics[width=0.96\linewidth]{assets/fig_local_success_heatmaps.pdf}
\caption{Rebuilt local 001 success-rate heatmaps.  Raw finite-slope Pandrosion deteriorates quickly with degree on the chosen interval; Steffensen is strong but fragile at $d=64$; cubic H3 remains successful on all tested starts.}
\label{fig:local-heatmaps}
\end{figure}

The rebuilt numbers match the qualitative structure of the original 001 benchmark.  The raw finite-slope primitive is the geometric base layer, but it is not a robust high-degree method by itself.  Steffensen acceleration largely fixes the local basin until the highest tested degree.  The cubic H3 correction is the strongest method in this local scalar test.

\section{New chart-scaling benchmark}
The second benchmark applies the Thales/Riemann layer to monomial equations $z^p=x$.  The degrees are
\[
p\in\{3,5,8,16,32,64\}.
\]
For each degree, the target grid contains 241 values with
\[
\log_{10}x\in[-120,120]+0.137,
\]
so the benchmark spans 240 decades and avoids exact coincidences with the scale palette.  The modes are:
\begin{itemize}[leftmargin=2.2em]
\item \textbf{unscaled}: use the direct chart for $x\ge1$ and the reciprocal chart for $x<1$, with $B=1$;
\item \textbf{$K=1$}: use the two-chart $10^\mathbb Z$ root-scale palette;
\item \textbf{$K=4$}: use the two-chart palette $10^{m+j/4}$, $j=0,1,2,3$.
\end{itemize}
Every normalized solve starts from $u_0=1$ and stops after at most 25 updates.  Success means relative error at most $10^{-12}$ in the returned root $z$.

\begin{table}[htbp]
\centering
\caption{Chart-scaling benchmark for cubic H3 over the 241-target grid.  Evaluation medians are computed only on successful runs.}
\small
\input{tables/table_chart_h3_summary.tex}
\label{tab:chart-h3}
\end{table}

\begin{figure}[htbp]
\centering
\includegraphics[width=0.96\linewidth]{assets/fig_chart_h3_success_evals.pdf}
\caption{Cubic H3 benchmark over 240 decades.  The unscaled chart succeeds only near favorable magnitudes.  The $K=1$ two-chart lattice solves all degrees through $p=32$ but is still too coarse for part of $p=64$.  The $K=4$ palette makes the normalized root interval small enough for full success in this test.}
\label{fig:chart-success}
\end{figure}

\begin{figure}[htbp]
\centering
\includegraphics[width=0.74\linewidth]{assets/fig_normalized_log_distribution.pdf}
\caption{Distribution of normalized $\log_{10}y$ values after chart selection.  Increasing the number of phase offsets shrinks the logarithmic covering radius before the finite-slope method begins.}
\label{fig:log-distribution}
\end{figure}

\begin{table}[htbp]
\centering
\caption{Selected H3 runs.  A dagger means the run did not meet the $10^{-12}$ returned-root tolerance.  The $K=4$ rows show the chosen chart and scale before the final return map.}
\small
\input{tables/table_chart_examples.tex}
\label{tab:examples}
\end{table}

The chart-scaling benchmark shows a different effect from the local polynomial benchmark.  In the local benchmark, all starts are already near the root.  In the monomial benchmark, the unscaled normalized target can be as large as $10^{120}$, so the same finite-slope formula forms probes in a badly scaled problem.  The Thales/Riemann layer changes the problem before the update is evaluated.  For $q=10$ and $K=4$, the normalized root is always in $[1,10^{1/8}]$, and cubic H3 succeeds on every tested target.

\section{Discussion}
The revised message is narrower and cleaner than a speed claim.  The finite-slope cubic update is not intended to replace exact-derivative Halley or Newton on ordinary monomial root extraction.  The point is that the Pandrosion program has two independent layers:
\begin{enumerate}[leftmargin=2.2em]
\item a derivative-free local correction ladder: raw finite slope, Steffensen, cubic H3;
\item a geometric preconditioner: Thales homothety plus optional reciprocal Riemann chart.
\end{enumerate}
The second layer is especially important because the finite-slope probes are scale-sensitive.  If the normalized root is near $1$, the raw predictor displacement is a useful proxy for the true local error, and the symmetric finite-slope Halley correction sees the cubic cancellation.  If the problem is left unscaled, the same formula may spend its probes measuring magnitude rather than curvature.

The scale-palette theorem also explains the difference between $K=1$ and $K=4$.  With $K=1$ and $q=10$, the worst root-coordinate normalized range is $[1,\sqrt{10}]$.  That is already magnitude-free, but it can still be wide for high-degree finite-slope probes.  With $K=4$, the worst range shrinks to $[1,10^{1/8}]$, which is small enough for full success in the present benchmark.

\section{Limitations}
This note remains deliberately restricted.
\begin{itemize}[leftmargin=2.2em]
\item The cubic theorem is scalar and local.
\item The chart-scaling theorem is monomial: it uses the two exact return maps for $z^p=x$.
\item The benchmark uses deterministic double-precision arithmetic.  It measures robustness of the proposed pipeline, not asymptotic wall-clock optimality.
\item The result is derivative-free, but not evaluation-free: H3 uses five function evaluations per update under the no-cache counting convention used here.
\end{itemize}

\section{Conclusion}
Paper 001 originally showed that the finite-slope/Pandrosion ladder can be pushed from raw first-order behavior to Steffensen-like quadratic behavior and then to a cubic Halley-like correction without exact derivatives.  This rebuilt version adds the chart-scaling layer from Paper 000.  For monomial targets, the solver should not attack $z^p=x$ in the original magnitude.  It should first choose a Thales scale and, when advantageous, the reciprocal Riemann chart, solve the normalized equation $u^p=y$ near $u=1$, and return by the appropriate chart map.

The rebuilt benchmark supports this combined view.  Locally, cubic finite-slope Halley remains the strongest method in the 001 polynomial test.  Globally over many decades of monomial targets, chart scaling converts the same cubic step from a magnitude-sensitive local method into a robust normalized solver on the tested grid.  The conceptual contribution is therefore the coupling: derivative-free cubic finite-slope correction after logarithmic chart normalization.

\section*{Reproducibility note}
The tables and figures in this manuscript were generated by the included Python script \path{generate_benchmarks.py}.  The run outputs include the raw local benchmark CSV, the chart benchmark CSV, summary CSV files, and the LaTeX table snippets used above.

\section*{Acknowledgments}
Generative AI disclosure.  OpenAI Codex/GPT-style tools were used during manuscript preparation for LaTeX editing, benchmark scripting, and editorial consistency checks.  The author reviewed the AI-assisted output and remains responsible for the accuracy, originality, integrity, and references of the manuscript.

\begin{thebibliography}{9}
\bibitem{halley1694}
E. Halley, ``A new, exact, and easy method of finding the roots of equations generally,'' \emph{Philosophical Transactions}, 1694.
\bibitem{aitken1926}
A. C. Aitken, ``On Bernoulli's numerical solution of algebraic equations,'' \emph{Proceedings of the Royal Society of Edinburgh} 46 (1926), 289--305.
\bibitem{steffensen1933}
J. F. Steffensen, ``Remarks on iteration,'' \emph{Skandinavisk Aktuarietidskrift} 16 (1933), 64--72.
\bibitem{traub1964}
J. F. Traub, \emph{Iterative Methods for the Solution of Equations}, Prentice-Hall, 1964.
\bibitem{householder1970}
A. S. Householder, \emph{The Numerical Treatment of a Single Nonlinear Equation}, McGraw-Hill, 1970.
\bibitem{milnor2006}
J. Milnor, \emph{Dynamics in One Complex Variable}, 3rd ed., Princeton University Press, 2006.
\bibitem{paper000}
I. Besevic, \emph{A Monomial Pandrosion Scheme: Contraction, Chart Scaling, and Steffensen Acceleration for $z^p=x$}, manuscript, 2026.
\bibitem{paper001old}
I. Besevic, \emph{Cubic Pandrosion by Finite-Slope Halley Acceleration}, manuscript, 2026.
\end{thebibliography}

\end{document}
